\documentclass[12pt]{article}
\usepackage[T1]{fontenc}
\usepackage[utf8]{inputenc}
\usepackage{lmodern}
\usepackage{textcomp}
\usepackage{enumitem}
\usepackage{geometry}
\geometry{margin=1in}
\begin{document}
\begin{center}
Answer Key - Constitutional Law - Undergraduate Level Assessment 8 \
\small (Answers are ordered exactly as in the question paper.)
\end{center}

\section*{Multiple Choice}
\begin{enumerate}[label=\arabic*.]
\item \textbf{Answer:} C
\\ \textit{Options:} A) There are fundamental differences between individual women.'; B) Men and women have different conceptions of the feminist project.'; C) Women look to context, whereas men appeal to neutral, abstract notions of justice.'; D) Men are unable to comprehend their differences from women.'
\\ \textit{Note:} Correct option: C - Women look to context, whereas men appeal to neutral, abstract notions of justice.'
\item \textbf{Answer:} C
\\ \textit{Options:} A) Parties to a treaty should be cognizant of its terms and not misinterpret them; B) Parties to a treaty should safeguard the object and purpose of the treaty; C) Parties to a treaty should adhere to its terms in good faith; D) Parties to a treaty should not violate the most important provisions of the treaty
\\ \textit{Note:} Correct option: C - Parties to a treaty should adhere to its terms in good faith
\item \textbf{Answer:} B
\\ \textit{Options:} A) It is a sin for humans not to apply reason.; B) The principles of natural law are discoverable by reason.; C) Natural law does not apply without good reason.; D) The law of nature is the basis of all positive law.
\\ \textit{Note:} Correct option: B - The principles of natural law are discoverable by reason.
\item \textbf{Answer:} A
\\ \textit{Options:} A) Article 3; B) Article 8; C) Article 9; D) Article 11
\\ \textit{Note:} Correct option: A - Article 3
\item \textbf{Answer:} C
\\ \textit{Options:} A) Treaties should always be designated as such and assume a particular form; B) Treaties should always assume a particular form, no matter how they are designated; C) Treaties do not have to assume a particular form or designated as such; D) Treaties have to be designated as such, no matter what form they assume
\\ \textit{Note:} Correct option: C - Treaties do not have to assume a particular form or designated as such
\end{enumerate}

\section*{True / False}
\begin{enumerate}[label=\arabic*.]
\setcounter{enumi}{5}
\item \textbf{Answer:} True
\\ \textit{Options:} A) True; B) False
\\ \textit{Note:} LLM-generated true statement based on MCQ.
\item \textbf{Answer:} True
\\ \textit{Options:} A) True; B) False
\\ \textit{Note:} LLM-generated true statement based on MCQ.
\item \textbf{Answer:} True
\\ \textit{Options:} A) True; B) False
\\ \textit{Note:} LLM-generated true statement based on MCQ.
\item \textbf{Answer:} True
\\ \textit{Options:} A) True; B) False
\\ \textit{Note:} LLM-generated true statement based on MCQ.
\item \textbf{Answer:} True
\\ \textit{Options:} A) True; B) False
\\ \textit{Note:} LLM-generated true statement based on MCQ.
\end{enumerate}

\section*{Long Form Response}
\begin{enumerate}[label=\arabic*.]
\setcounter{enumi}{10}
\item \textbf{Answer:} It illuminates the concept of a rule.. Explain why this is the correct answer and provide supporting evidence. Reference key facts or concepts that support this choice.
\\ \textit{Note:} Long-form derived from MCQ; award credit for stating the correct idea and giving supporting reasoning.
\item \textbf{Answer:} The international recognition of human rights after World War II. Explain why this is the correct answer and provide supporting evidence. Reference key facts or concepts that support this choice.
\\ \textit{Note:} Long-form derived from MCQ; award credit for stating the correct idea and giving supporting reasoning.
\end{enumerate}

\vfill
\noindent\textit{End of Answer Key}
\end{document}